\documentclass{article}
\title{{Practice Problem 2.38}}
\author{CSAPP}
\usepackage{booktabs, parskip }

\begin{document}
    \maketitle
    
    \begin{flushleft}
        The $LEA$ instruction can perform computations of the form
        $(a << k) + b$ where $k$ is either $0,1,2$ or  $3$ and $b$ is either $0$ or some
        program value. The compiler often uses this instruction to perform multiplication
        by constant factors. For example, we can compute $3 * a$ as $(a << 1) + a$
        Considering the cases where $b$ is either $0$ or equal to $a$, all possible 
        values of $k$, what multiples of $a$ can be computed with a single $LEA$ instruction?

        \smallskip

        \textbf{\small SOLUTION}
        For each value of $k$, we can compute two multiplies: $2^k$
        (where $b$ is $0$) and $2^k + 1$  (when $b$ is $a$).
        Thus, we can compute multiples $1,2,3,4,5,8$ and $9$
        
    \end{flushleft}
\end{document}